\section{What a replacement validation would need to establish}
\label{sec:validation}

The case identifies three decisions that must precede another cutpoint claim.
First, fix the inferential target. A conditional cell-sampling study can hold a
source pool fixed and evaluate intervals against $f_N(s-1)\mu_{\mathrm{pool}}$.
A patient-population study instead needs a patient-sampling model and a target
averaged over that population. The current draw-dependent benchmark answers a
third question. Results under one definition cannot be relabelled as coverage
under another. The same distinction applies to a clinical simulator: requested
population treatment effects, latent finite-cohort effects and empirical
observed-data contrasts are separate validation targets.

Second, assess the whole decision rule using the interval that would actually
be reported. Record bias and error against the chosen target, null rejection,
detection at the specified effect, and abstention frequency. Report both
conditional interval performance among returned answers and how often a valid
answer is returned. Our five-cell diagnostic concerns arithmetic before a gate
that already abstains there; the 50-cell null results concern its reporting
boundary directly. A successful effect-detection rate can accompany excessive
null rejection, so the former alone cannot validate an interval. A fixed
5\% null target is appropriate here because the reported interval is nominally
95\%; other operating targets must follow the intended use.

Third, separate selecting a rule from assessing it. If a scalar threshold is
intended to guarantee performance above that threshold, qualification must
address the evaluated upper range rather than only its first passing point.
Reserve independent simulation draws for the final assessment, report rate
uncertainty and the tested range, and distinguish this from transfer to new
patients or cohorts. Increasing resamples from an empirical pool improves
conditional numerical precision; it does not broaden that pool's biological
support. For an external-control application, the additional question is
comparability across data sources, which requires direct study-level validation
such as the held-out comparisons of \citet{ventz2019}. The present experiments
motivate these requirements. The next paragraph performs the first of them.

\paragraph{A constructed external-control workflow where the first requirement
decides the answer.} The requirement above is that fixing the inferential
target must precede everything else. We built a workflow in which it does, and
in which getting it wrong is not caught by any routine safeguard
(Appendix~\ref{app:eca}; pre-registered with nine falsifiers, none fired).

A single-arm trial draws an external control arm three times its size from a
disease registry, and is validated by plasmode resampling beforehand.
Enrolment is without replacement under eligibility weights enriching for
refractory patients; the induced benefit is effect-modified and each record
carries a latent responsiveness, so \emph{no scalar in the configuration equals
the marginal truth}. This is the setting the audit is for: a requested
parameter is genuinely unavailable, and a reference must be computed from the
sample.

The injected defect is a single misidentified target population --- the
estimator answers for the pooled analysis file rather than the enrolled
patients, reachable either as an inverse-probability weight written $1/e$
instead of $e/(1-e)$ or as averaging stratum contrasts over the analysis file.
Because the control arm is three times the trial, the pooled mix is the
registry's.

\textbf{The validation does not merely pass; it ranks the defective estimator
first.} Against the observed-data reference, standardisation for the pooled
population has maximum residual $0$ bitwise at every cohort size, recovery
ratio $1.000000$, and the lowest error of the seven estimators, while the
correct enrolled-population estimator appears $46\%$ attenuated. A bake-off
run this way instructs the analyst to discard the correct estimator and deploy
the broken one. Positivity holds and the post-weighting balance table is
immaculate --- worst between-arm standardised mean difference
$5.6\times10^{-16}$ --- while the weighted pseudo-population sits $0.964$ from
the patients actually enrolled, a column no protocol tabulates. The reported
benefit is $1.82$--$1.93$ months against a pre-specified clinically important
difference of $1.5$; the enrolled patients received $0.97$--$1.00$. The
overstatement is $1.89$-fold and it crosses the decision threshold.

\textbf{Two pre-registered controls corrected our own prescription.} Replacing
the reference with realised individual effects among the enrolled inverts the
ranking completely: the correct estimator becomes unbiased with error falling
as $n^{-1/2}$, and the defective one carries a flat $+0.885$. But a reference
that merely breaks the equality is \emph{not} sufficient --- the same
potential-outcome reference computed on the pooled population has a nonzero
residual and still hides the defect, at bias $+0.003$. \textbf{A reference must
name the target population, not merely differ from the estimator}, which is
stronger than the prescription we gave in \S\ref{sec:blind}. And at zero
case-mix shift the equality persists while the harm disappears: equality marks
a validation that could not have failed, never a wrong answer.

This workflow is constructed, not observed. It shows the failure is reachable
by ordinary choices and expensive when reached; it is not evidence that anyone
has made it, and the repository search in Appendix~\ref{app:prevalence} remains
negative.
